\documentclass[a4paper]{article}

\input{preamble.sty}
\input{summary-preamble.sty}

\title{Crash Course i AA MM5020}
\author{Disa Degerholm}
\date{\today}

\lhead{MM5020}
\chead{\textbf{Formelblad - Abstrakt Algebra}}
\rhead{Disa Degerholm | \today}

\begin{document}
\makemytitlepage

\section{Groups}
\subsection{Groups and subgroups}
\begin{defi}{Group properties \& algebraic structures}
Let $G \neq \emptyset$ be a set with a map $\circ : G \times G \to G$. We define the following properties (with the implicit operation $\circ$):
\begin{itemize}
    \item \textbf{Associative:} $\forall a,b,c \in G : (ab)c = a(bc)$ 
    \item \textbf{Unit element/Identity:} $\exists e \in G : ea = a = ae$ 
    \item \textbf{Inverse elements:} $\forall a \in G, \exists b \in G : ab = e = ba$ 
    \item \textbf{Commutative:} $\forall a,b \in G : ab = ba$ 
\end{itemize}
Then, $(G, \circ)$ is called a 
\begin{itemize}
    \item \textbf{Semigroup} if associativity holds.
    \item \textbf{Monoid} if associativity and existence of a unit hold.
    \item \textbf{Group} if associativity, existence of a unit, and existence of inverses hold.
    \item \textbf{Abelian group} if associativity, existence of a unit, existence of inverses, and commutativity hold.
\end{itemize}
\end{defi}

\begin{ther}{Uniqueness of unit and inverse}
Let $G$ be a group then the following holds:
\begin{enumerate}[label=(\alph*)]
    \item The unit element is unique.
    \item The inverse of any $a \in G$ is unique.
\end{enumerate}
\end{ther}
\begin{add}
    Proof exists in KTH lecture notes.
\end{add}

\begin{ther}{Weak criteria for a group}
Let $G \neq \emptyset$ be a set with a map $\circ : G \times G \to G$ such that $\circ$ is associative:
\begin{itemize}
    \item There exists a left unit: $\exists e \in G, \forall a \in G : ea = a$.
    \item All elements have a left inverse: $\forall a \in G, \exists b \in G : ba = e$.
\end{itemize}
Then $(G, \circ)$ is a group.
\end{ther}
\begin{add}
    Proof exists in KTH lecture notes.
\end{add}

\begin{defi}{Submonoid}
Let $G$ be a monoid and $H \subseteq G$ a subset. $H$ is a submonoid of $G$ if 
\begin{enumerate}[label=(\roman*)]
    \item $e \in H$
    \item $\forall a,b \in H : ab \in H$
\end{enumerate}
\end{defi}

\begin{defi}{Subgroup}
Let $G$ be a group and $H \subseteq G$ a subset. $H$ is a subgroup of $G$, written $H \le G$, if 
\begin{enumerate}[label=(\roman*)]
    \item $e \in H$ 
    \item $\forall a,b \in H : ab \in H$ 
    \item $\forall a \in H : a^{-1} \in H$ 
\end{enumerate}
\end{defi}
\begin{add}
    A trivial subgroup is the always exisiting subgroup $\{e\}$ alone or together with the rest of the group, since $G\leq G$. While a non-trivial subgroup is any subgroup $H$ such that $G\neq H \neq \{e\}$ i.e the same identity element but partly the same or other extra elements.
\end{add}

\subsection{Group homomorphisms}

\begin{aha}
    Let $(G, \circ)$ and $(H, \cdot)$ be two groups where $\cdot$ and $\circ$ are groupoperations. A function $f : G \to H$ is a \key{group homomorphism} if $\forall a,b \in G : f(ab) = f(a)f(b)$ (can be written as $G\sim H$). In other words, the function $f$ preserves the group operation. 

    Meanwhile \key{group isomorphisms} are homomorphisms (same notation as above) that are bijective, i.e where no information is destroyed or created. It is written as $G \cong H$ if at least one isomorphism exists between the groups. In other words, the function $f_2$ preserves the group operation and is a bijection. In other words isomorphisms$\subseteq$ homomorphims but are NOT usually subgroups. 
    
    A simple example of homomorphism is the exponential function $f : (\mathbb{R}, +) \to (\mathbb{R}^+, \cdot)$, where $f(x+y) = e^{x+y} = e^x e^y = f(x)f(y)$ but it is not an isomorphism since it is not injective (since $e^0 = e^{2\pi i}$ for example). However the function $f : (\mathbb{R}, +) \to (\mathbb{R}, +)$, $x \mapsto 2x$ is an isomorphism since it is a homomorphism and bijective by the fact that $2x = 2y \implies x=y$ and $\forall y \in \mathbb{R}, \exists x \in \mathbb{R} : 2x = y$.
\end{aha}

\begin{defi}{Monoid homomorphism}
Let $G_1, G_2$ be monoids and $\varphi : G_1 \to G_2$ a map. $\varphi$ is a monoid homomorphism if
\begin{enumerate}[label=(\roman*)]
    \item $\forall a,b \in G_1 : \varphi(ab) = \varphi(a)\varphi(b)$
    \item $\varphi(e_{G_1}) = e_{G_2}$
\end{enumerate}
\end{defi}

\begin{defi}{Group homomorphism}
Let $G_1, G_2$ be groups and $\varphi : G_1 \to G_2$ a map. $\varphi$ is a group homomorphism if
\begin{enumerate}[label=(\roman*)]
    \item $\forall a,b \in G_1 : \varphi(ab) = \varphi(a)\varphi(b)$
    \item $\varphi(e_{G_1}) = e_{G_2}$
\end{enumerate}
\end{defi}

\begin{ther}{Properties of group homomorphisms}
Let $\varphi : G_1 \to G_2$ be a group homomorphism. Then:
\begin{enumerate}[label=(\alph*)]
    \item $\varphi(e_{G_1}) = e_{G_2}$
    \item $\forall a \in G_1 : \varphi(a)^{-1} = \varphi(a^{-1})$
\end{enumerate}
\end{ther}
\begin{add}
    This also holds true if let $\varphi : G_1\to G_2$ be a map between groups and $\varphi(ab)=\varphi(a)\varphi(b)$ $\forall a, b \in G_1$ then $\phi$ is a group isomorphism and property b holds. OBS proof exists in KTH lecture notes. 
\end{add}

\begin{defi}{Types of group homomorphisms}
\begin{enumerate}[label=(\alph*)]
    \item For two groups $G_1, G_2$: $\text{Hom}(G_1,G_2) := \{\varphi : G_1 \to G_2 \mid \varphi \text{ is a group homomorphism}\}$.
    \item An injective homomorphism is a monomorphism. A surjective homomorphism is an epimorphism. A bijective homomorphism is an isomorphism.
    \item For a group $G$: $\text{End}(G) := \text{Hom}(G,G)$ and its elements are endomorphisms of $G$ (i.e Hom(Z, Z), where Z is arbitrary). Also, $\text{Aut}(G) := \{\varphi \in \text{End}(G) \mid \varphi \text{ is bijective}\}$ and its elements are automorphisms of $G$.
\end{enumerate}
\end{defi}

\begin{ther}{Composition of homomorphisms}
Compositions of group homomorphisms are group homomorphisms. Ex $g_1, G_2, G_3$ are groups and $\varphi_1 \in \text{Hom}(G_1, G_2)$, $\varphi_2 \in \text{Hom}(G_2, G_3)$ then $\varphi_2 \circ \varphi_1 \in \text{Hom}(G_1, G_3)$.
\end{ther}
\begin{add}
    Proof exists in KTH lecture notes.
\end{add}

\begin{ther}{Automorphism group}
If $G$ is a group, then $(\text{Aut}(G), \circ)$ is a group.
\end{ther}
\begin{add}
    Proof exists in KTH lecture notes.
\end{add}

\begin{defi}{Kernel and image of a group homomorphism}
\begin{enumerate}[label=(\alph*)]
    \item The kernel of a group homomorphism $\varphi : G_1 \to G_2$ is $\text{ker}\,\varphi := \{a \in G_1 \mid \varphi(a) = e_{G_2}\}$.
    \item The image of a group homomorphism $\varphi : G_1 \to G_2$ is $\text{im}\,\varphi := \{\varphi(a) \mid a \in G_1\}$.
\end{enumerate}
\end{defi}

\begin{ther}{Kernel, image and subgroups under homomorphisms}
Let $G_1$ and $G_2$ be groups and $\varphi \in \text{Hom}(G_1, G_2)$.
\begin{enumerate}[label=(\alph*)]
    \item $\text{ker}\,\varphi \le G_1$
    \item $\text{im}\,\varphi \le G_2$
    \item $\varphi \text{ is injective} \iff \text{ker}\,\varphi = \{e\}$
    \item $\varphi \text{ is surjective} \iff \text{im}\,\varphi = G_2$
    \item $H_1 \le G_1 \implies \varphi(H_1) \le G_2$
    \item $H_2 \le G_2 \implies \varphi^{-1}(H_2) \le G_1$
\end{enumerate}
\end{ther}
\begin{aha}
    OBS $\varphi(e_1)=e_2$. 
\end{aha}
\begin{add}
    Proof exists in KTH lecture notes.
\end{add}

\begin{defi}{Conjugation}
Let $G$ be a group and $a \in G$ an element. We define conjugation by $a$ as $\gamma_a : G \to G$, $g \mapsto aga^{-1}$; also called an inner automorphism of $G$. 
$$\text{Inn}(G) := \{\gamma_a \mid a \in G\} \le \text{Aut}(G)$$
\end{defi}

\begin{add}
It holds that $\gamma_a \in \text{Aut}(G)$ and the function that creates all $\gamma_a$ i.e $\phi: G \to \text{Aut}(G)$ where $a \mapsto \gamma_a$ is a group homomorphism (in other words $\gamma_{ab} = \gamma_a \circ \gamma_b$).
\end{add}

\subsection{Cosets}

\begin{defi}{Coset}
Let $G$ be a group and $H \le G$.
\begin{enumerate}[label=(\alph*)]
    \item A \textbf{left coset} of $H$ in $G$ is a subset of the form 
    $$aH := \{ah \mid h \in H\}, \quad a \in G$$
    \item $G/H := \{aH \mid a \in G\}$ denotes the set of all left cosets of $H$ in $G$.
    \item Analogously, $Ha := \{ha \mid h \in H\}$ (for $a \in G$) is a \textbf{right coset} of $H$ in $G$. The set of all right cosets is denoted $H \backslash G := \{Ha \mid a \in G\}$.
\end{enumerate}
\end{defi}

\begin{add}
The bijection $G \to G$, $g \mapsto g^{-1}$ restricts to a bijection $aH \to Ha^{-1}$. Hence, it defines a bijection $G/H \to H \backslash G$, $aH \mapsto Ha^{-1}$.
\end{add}

\begin{ther}{Equivalent coset conditions}
Let $G$ be a group, $H \le G$, and $a,b \in G$. Then the following are equivalent:
\begin{enumerate}[label=(\roman*)]
    \item $aH = bH$,
    \item $aH \cap bH \neq \emptyset$,
    \item $a \in bH$,
    \item $b^{-1}a \in H$.
\end{enumerate}
\end{ther}

\begin{ther}{Partition of $G$}
Let $H \le G$. Then $G$ is the disjoint union (union where they share nothing) of all left cosets of $H$ in $G$, i.e.,
$$G = \dot{\bigcup_{C \in G/H}} C$$
\end{ther}

\begin{defi}{Coset properties}
Let $G$ be a group, $H \le G$:
\begin{enumerate}[label=(\alph*)]
    \item The elements of a coset $aH$ in $G$ are called the \textbf{representatives of the coset}. 
    \item $[G : H] := |G/H| = |H \backslash G|$ is the \textbf{index of $H$ in $G$}. Intuativly, how many times G is bigger than H. 
    \item $\text{ord}(G) := |G|$ is the \textbf{order of $G$} (nr of elements in G).
\end{enumerate}
\end{defi}

\begin{ther}{Size of cosets}
Let $G$ be a group, $H \le G$, and $a,b \in G$. There is a bijection $aH \to bH$. In particular, $|aH| = |bH| = \text{ord}(H)$ (because h$\mapsto$ah is bijective)a.
\end{ther}

\begin{ther}{Theorem of Lagrange}
Let $G$ be a finite group, $H \le G$. Then:
$$\text{ord}(G) = \text{ord}(H) \cdot [G : H]$$
\end{ther}

\subsection*{Normal subgroups}

\begin{defi}{Normal subgroup}
Let $G$ be a group, $H \le G$. $H$ is a \textbf{normal subgroup} of $G$ (intuition balanced/symetric subgroup like 0 or 1 to create group division and kernal of homomorphisms), denoted $H \trianglelefteq G$, if:
$$\forall \gamma \in \text{Inn}(G) : \gamma(H) = H, \quad \text{i.e}, \forall a \in G : aHa^{-1} = H \text{. Ex G=(Z, +) and H=(2Z, +)}.$$
\end{defi}

\begin{ther}{Normal subgroup criteria}
Let $G$ be a group, $H \le G$. Then the following are equivalent:
\begin{enumerate}[label=(\roman*)]
    \item $H \trianglelefteq G$,
    \item $\forall \gamma \in \text{Inn}(G) : \gamma(H) \subseteq H$,
    \item $\forall \gamma \in \text{Inn}(G) : H \subseteq \gamma(H)$,
    \item $\forall a \in G : aH = Ha$.
\end{enumerate}
\end{ther}

\begin{add}
In this case, $aH = Ha$ is called the \textbf{residue class of $a$ modulo $H$} (left coset to H under a is equal to the right).
\end{add}

\begin{ther}{Quotient groups}
Let $G$ be a group, $N \trianglelefteq G$:
\begin{enumerate}[label=(\alph*)]
    \item $\forall a,b \in G, \ (aN)(bN) = abN$ (where N is not a siingel element but the whole normal subgroup).
    \item $G/N={aN \mid a \in G}$ is a group with multiplication as in (a), called the \textbf{factor group} or \textbf{quotient group} of $G$ modulo $N$. Its unit element is $eN = N$. (Inutition: work with N instead of e and you get a group with the same structure as G but with N as the new identity element (N becomes e) and the rest of the elements are the same but with N added to them (aN instead of a, like 1*nr)).
    \item The \textbf{canonical projection} $\pi : G \to G/N, \ a \mapsto aN$, is a surjective group homomorphism with $\text{ker}(\pi) = N$.
\end{enumerate}
\end{ther}

\begin{ther}{Universal property of $\pi$}
Let $\varphi : G_1 \to G_2$ be a group homomorphism, and let $N \trianglelefteq G_1$ such that $N \subseteq \text{ker}(\varphi)$. Then there is a unique group homomorphism $\bar{\varphi} : G_1/N \to G_2$, such that $\varphi = \bar{\varphi} \circ \pi$, i.e., the following diagram is commutative:

\includegraphics[width=0.5\textwidth]{Bilder\AA_universal_property.png}

Moreover:
\begin{enumerate}[label=(\alph*)]
    \item $\text{im}(\varphi) = \text{im}(\bar{\varphi})$,
    \item $\text{ker}(\bar{\varphi}) = \pi(\text{ker}(\varphi))$,
    \item $\text{ker}(\varphi) = \pi^{-1}(\text{ker}(\bar{\varphi}))$,
    \item $\bar{\varphi}\text{ is injective} \iff N = \text{ker}(\varphi)$.
\end{enumerate}
\end{ther}

\begin{ther}{First isomorphism theorem}
If $\varphi : G_1 \to G_2$ is a surjective group homomorphism, then $G_2$ is canonically isomorphic to $G_1/\text{ker}(\varphi)$.
\end{ther}

\subsection{Cyclic groups}
basdjfadbfvasuofbvuiasbfjkasbfjöasbfjkö hcuaophf aOHP DIAJO CNAO NAKO CNAKJLÖS

\subsection{Group actions}
basdjfadbfvasuofbvuiasbfjkasbfjöasbfjkö hcuaophf aOHP DIAJO CNAO NAKO CNAKJLÖS

\subsection{Sylow groups}
basdjfadbfvasuofbvuiasbfjkasbfjöasbfjkö hcuaophf aOHP DIAJO CNAO NAKO CNAKJLÖS

\subsection{Permutation groups}
basdjfadbfvasuofbvuiasbfjkasbfjöasbfjkö hcuaophf aOHP DIAJO CNAO NAKO CNAKJLÖS

\subsection{Some notation for important groups}
basdjfadbfvasuofbvuiasbfjkasbfjöasbfjkö hcuaophf aOHP DIAJO CNAO NAKO CNAKJLÖS

\section{Rings}
\subsection{Rings, subrings \& ideals}
\subsection{Ring homomorphisms}
\subsection{Polynomials}
\subsection{Integral domains}
\subsection{Primes}
\subsection{Euclidian domains}
\subsection{Unique factorization domains}
\subsection{Modules}
\subsection{Field extensions}
\subsection{Algebraic closures}
\subsection{Splitting fields}
\subsection{Finite fields}
\subsection{Overview of types of rings}

\section{Other good math tricks}

Bijectivity and injectivity means that f(x)=f(y) -> x=y


\section{Kommandon}

\begin{defi}{Vandermondmatris}
    Här är def 2
\end{defi}

\begin{ther}{banan satsen}
    Här är theorom nr 2 med ett $\key{nyckelord}$
\end{ther}

\begin{prf}{linjäriseringen}
Här är proof nr2
\end{prf}

\begin{add}
Här är kommentar nr2
\end{add}

\begin{aha}
Här är förståelse nr2
\end{aha}

\end{document}